%!TEX root = ../template.tex
%%%%%%%%%%%%%%%%%%%%%%%%%%%%%%%%%%%%%%%%%%%%%%%%%%%%%%%%%%%%%%%%%%%%
%% chapter8.tex
%% NOVA thesis document file
%%
%% Chapter with lots of dummy text
%%%%%%%%%%%%%%%%%%%%%%%%%%%%%%%%%%%%%%%%%%%%%%%%%%%%%%%%%%%%%%%%%%%%

\typeout{NT FILE chapter8.tex}%

\chapter{Evaluation and future work}
\label{cha:Evaluation}


\section{Evaluation}
\label{sec:evaluation}

During the implementation of the main algorithms of the OCaml-FLAT library, unit tests were written for all operations to provide greater assurance of 
their correct functioning. Especially on tests for calculating attributes, since 
the calculation of inherited attributes was more complex than expected, and it is one of the main focuses of this dissertation. 
Care was also taken to design unit tests that exercise potentially problematic cases, for example attribute grammars that create cyclic situations by 
continuously changing attribute values. Nevertheless, it would have been preferable to make the application available to a limited number of users so it 
could be tested more intensively. To help evaluate the quality and correctness of the work, the supervisor of this dissertation was frequently consulted 
in his capacity as a former lecturer of the Theory of Computation course.

\section{Conclusion}
\label{sec:conclusion}

For this project the OCaml-FLAT library was extended; the OFLAT graphical tool exists to create and manipulate FLAT models, but I did not work on 
the OFLAT library.

The OCaml-FLAT library was adapted to support attribute grammars, with all necessary types created for their representation and a range of operations 
implemented. These include evaluation of attribute values for a given input, generation of words of length n, conversions between the various types 
implemented in the library, and transformations between different grammar models. In addition, functions were implemented to produce derivation trees 
that represent the history of a derivation or evaluation.


\section{Future Work}

Future work includes developing the graphical front end and extending support for attribute grammars in the LL(1) and LR classes, 
since OFLAT already supports these grammar families and they have the advantage of defining deterministic parsers (for example, tools such as Yacc).  

A good next step would also be to calculate attributes for LL(1) grammars, this would 
require implementing a dedicated LL(1) parser module that interleaves parsing with attribute evaluation. In practice, attribute values would be computed at 
runtime by the deterministic parser, taking care to handle both synthesized and inherited attributes and to detect or prevent cyclic attribute dependencies.  

Additional directions include integrating the graphical interface with the parser modules so users can build, validate and execute attribute grammars 
interactively, and adding diagnostics and visualization tools to help identify problematic cases.
